% ECONOMICS_PROBLEM: {"id": "H24-183", "title": "Inheritance taxation and gifts", "jel_code": "H24", "source_id": "OP-0473"}
% !TeX program = xelatex
\documentclass[11pt,a4paper]{article}
\usepackage[margin=23mm]{geometry}
\usepackage{fontspec}
\setmainfont{Latin Modern Roman}
\usepackage{amsmath,amssymb,mathtools}
\usepackage{xcolor,enumitem,booktabs,longtable,array,xurl,hyperref}
\definecolor{muted}{HTML}{53636C}
\hypersetup{colorlinks=true,urlcolor=blue,linkcolor=blue}
\setlength{\parindent}{0pt}
\setlength{\parskip}{5pt}
\newcommand{\R}{\mathbb R}
\newcommand{\E}{\mathbb E}
\newcommand{\N}{\mathbb N}
\newcommand{\ind}{\mathbf 1}
\newcommand{\Var}{\operatorname{Var}}
\newcommand{\Cov}{\operatorname{Cov}}
\newcommand{\supp}{\operatorname{supp}}
\DeclareMathOperator*{\argmin}{arg\,min}
\DeclareMathOperator*{\argmax}{arg\,max}
\newcommand{\entrymeta}[3]{{\sffamily\small\color{muted}\textbf{\texttt{#1}}\quad|\quad #2\par #3\par}}
\newcommand{\Problemref}[1]{\texttt{#1}}
\newcommand{\JELref}[2]{\texttt{#2} (JEL \texttt{#1})}
\begin{document}
\section*{Inheritance taxation and gifts}
\textbf{Problem ID:} \texttt{H24-183}\quad \textbf{JEL:} \texttt{H24}\par
% BEGIN ECONOMICS STATEMENT
\label{problem:OP-0473}
{\sffamily\small\color{muted}Original subject group: Taxation and social insurance\par}
\label{definitions:problem:30007755}
\textbf{Audit classification:} Published research agenda. IFS section 7.5 (lifetime gifts), Box 7.5, and section 7.6 support a research agenda. Corrected conditioning on donor death and explained that the regression objective is editorial and not a causal or comprehensive mobility measure.
\subsection*{Definitions and precise research target}
Fix a cohort of donor--recipient families before reform, finite horizon \(H\), fixed numeraire and parental baseline resources \(P_i\) with \(0<\operatorname{Var}(P_i)<\infty\). A regime \(p\in\mathcal P\) specifies estate and lifetime-gift taxes, exemptions, reporting, enforcement, public-revenue use and the complete anticipation information. Death times are random; all families are followed through \(H\), including those whose donor survives. Define gifts before \(H\), estates transferred before \(H\) (zero if no death), unreported transfers and saving under each regime. Let \(Y_i^H(p)\) be recipient terminal resources, including all net transfers received by \(H\); do not condition on policy-dependent death, gift receipt or residence. Require second moments and define
\[
 \beta(p)=\frac{\operatorname{Cov}(Y_i^H(p),P_i)}{\operatorname{Var}(P_i)}.
\]
This descriptive association is not the causal effect of parental resources and is sensitive to the units of \(P_i\) and \(Y_i\). Fix a cardinal family welfare index \(V_i(p)\), baseline \(p_0\), welfare tolerance \(\delta\geq0\), and required discounted receipts \(\bar R\). The editorial optimization target is
\[
 \inf_{p\in\mathcal P}\beta(p),\qquad
 R(p)\geq\bar R,\quad \mathbb E[V_i(p)-V_i(p_0)]\geq-\delta.
\]
Report feasibility and optimizer attainment separately. A structural model must specify donor motives, recipient behavior, concealed transfers, gift timing, mortality and equilibrium selection. With data law \(P_O\), derive the joint set of \((\beta(p),R(p),\mathbb EV_i(p))_{p\in\mathcal P}\) over compatible models; its coupling across policies matters for robust design. The sign of a regression slope is not a complete inequality criterion: minimizing it can reward reversals of parental advantage. A monotonicity restriction or a separate inequality objective is needed if such reversals are excluded. The IFS source motivates behavioral and gift-design research, not this regression-based policy criterion.
\subsection*{Variants and assumption changes}
\begin{enumerate}
\item \textbf{Estate tax versus recipient accession tax (editorial).} Tax donor estates in one class; tax each recipient's cumulative receipts in another. Specify lifetime record linkage and equivalent revenue constraints before comparing objectives.
\item \textbf{Linear association versus rank mobility (editorial).} Replace \(\beta\) by a slope of population midranks, defining ties by mid-distribution ranks. This removes sensitivity to monetary scale but changes the welfare target.
\item \textbf{Fixed versus endogenous gifts (editorial).} Freeze gift timing as a mechanical benchmark, then allow timing and asset-form responses. A static costing is a different counterfactual from behavioral reform effects.
\end{enumerate}
\subsection*{References and source audit}
Advani--Sturrock (2023), IFS report: section 7.5, Treatment of gifts; Box 7.5; section 7.6. Metadata and explicit further-work passages are verified. The slope-constrained optimization is newly authored.
\begin{enumerate}\raggedright
\item\hypertarget{definitions:source:30007755:1}{} Arun Advani; David Sturrock. 2023. \emph{Reforming inheritance tax}. Section7.5, Treatment of gifts; Box7.5, Responses to tax changes; section7.6, Effects on social mobility.
\url{https://ifs.org.uk/publications/reforming-inheritance-tax}.
DOI: \url{https://doi.org/10.1920/re.ifs.2023.0275}.
\end{enumerate}

\subsection*{Common conventions: Source mathematical conventions}

These conventions apply to each entry unless explicitly replaced.

\textbf{Models and identification.} A primitive model \(m\) specifies all probability laws, feasible choices, information, equilibrium and measurement rules used by its target. The admissible class \(\mathcal M\) is fixed independently of outcomes. If \(P_O(m)\) is its observable probability law and \(\tau(m)\) the target, its identified set at observable law \(P\) is
\[
 \mathcal I_\tau(P)=\{\tau(m):m\in\mathcal M,\ P_O(m)=P\}.
\]
Point identification means that this set is a singleton. An empty set signals incompatibility of data and assumptions. Coordinate bounds need not describe the joint set. Bounds are sharp only if every retained target is attainable or arbitrarily approachable by compatible models. Finite bounds require explicit restrictions on unobserved quantities; unrestricted confounding, hidden wealth, extrapolation or arbitrary equilibrium selection can yield uninformative sets.

\textbf{Causal interventions.} A potential outcome \(Y_i(p)\) is the outcome under a complete policy regime \(p\), including timing, financing, information and market spillovers. The target population and units are fixed before treatment. With interference, use the complete assignment vector or a justified exposure mapping; individual no-interference notation alone cannot identify an economy-wide intervention. A difference of mean potential outcomes does not require the cross-world joint distribution. Conditioning on post-treatment migration, survival, enrollment or employment changes the estimand and needs separate justification.

\textbf{Statistical statements.} An estimator, confidence set, forecast or test specifies a sampling law, model class and loss/coverage criterion. Uniform \((1-\alpha)\)-coverage means \(\inf_{m\in\mathcal M}\Pr_m\{\tau(m)\in C_n\}\geq1-\alpha\), or an explicitly asymptotic version. The whole parameter space trivially has coverage, so informativeness requires a stated width, risk or power objective. A forecasting improvement is relative to a named benchmark, information set and evaluation population.

\textbf{Equilibrium and optimization.} An equilibrium comprises feasible plans, optimality, beliefs where needed, budget constraints and all market/resource clearing. The primitives or a stated benchmark define each correspondence \(\mathcal E(p;m)\). Nonemptiness is assumed or proved, never inferred just from compact policy sets. A selected-equilibrium target uses a declared selection rule; a robust target quantifies over all permitted equilibria. Use suprema or infima unless attainment is established. An \(\varepsilon\)-optimal policy is feasible and within \(\varepsilon\) of the finite optimum value. Report an empty feasible set separately.

\textbf{Welfare and accounting.} Normative utility, interpersonal normalization, social weights, numeraire, discounting and the use of public receipts are primitives. Behavioral utility need not coincide with the welfare criterion. Household welfare includes ownership income and rents once; adding profits again to compensating variation generally double-counts them. A monetary equivalent is defined by an explicit transfer and price convention, with monotonicity and range conditions. Average effects, mechanism contributions, transfers and welfare losses are different targets. Infinite sums require convergence and dynamic feasible plans require terminal or no-Ponzi conditions.

\textbf{Source versions.} A working-paper date, revision date and journal publication year are distinguished. A DOI for a journal version is not automatically the DOI of an earlier manuscript. Locators refer to the stated version and printed pages where specified. A metadata-only check does not verify a claim about a full-text conclusion. Every entry's source subsection narrows the provenance claim accordingly.
% END ECONOMICS STATEMENT
\end{document}
